\documentclass[10pt,a4paper]{amsart}
\usepackage[margin=19mm]{geometry}
\usepackage{amsmath,amssymb,mathtools}
\usepackage[scr=rsfs,cal=euler]{mathalfa}
\usepackage{xfrac}
\usepackage[unicode,hidelinks,bookmarks=false]{hyperref}
\newcommand{\ie}{i.e.\ }
\newcommand{\cf}{cf.\ }
\newcommand{\resp}{respectively\ }
\newcommand{\Subsection}{\subsection}
\providecommand{\nobreakdash}{\nobreak\hbox{-}}
\setlength{\emergencystretch}{2em}
\multlinegap0pt
\usepackage{bm}
\usepackage{bbm}
\usepackage{dsfont}
\usepackage{xspace}
\usepackage{cancel}
\usepackage{mathtools}
\usepackage{amsthm}
\makeatletter
\@ifundefined{thm}{\newtheorem{thm}{Theorem}}{}
\makeatother
\def\Tr{{\rm Tr}}
\newcommand{\F}{\mathbb F}
\title[Permutation polynomials from the trace functions]{Permutation polynomials from the trace functions}
\author{Sartaj Ul Hasan, Ramandeep Kaur, Hridesh Kumar}
\date{Source: arXiv:2608.10776v1 (CC BY 4.0)}
\begin{document}
\maketitle

\noindent\emph{Source and licence.} Permutation polynomials from the trace functions. Version arXiv:2608.10776v1 (CC BY 4.0), distributed under the Creative Commons Attribution 4.0 International licence, as recorded in the arXiv licence metadata for this exact version and shown on the abstract page https://arxiv.org/abs/2608.10776v1. Full rights notice: licenses/arxiv-2608.10776-CC-BY-4.0.txt.\par\smallskip

\noindent\textit{Editorial scope.} Counted unit: the Thm whose source label is \texttt{T31} in the article (source0.tex lines 246--249, source label \texttt{T31}), delivered together with the article's own complete original proof of it (source0.tex lines 250--320, one \texttt{proof} environment of 71 lines). No mathematics is written, repaired or re-derived: statement and proof are the article's own text line for line. Required context retained uncounted: none. Every definition, display and label of those lines is retained unchanged. Omitted from this package and not redistributed: 1--245, 321--769. Nothing inside a retained excerpt is omitted or altered: every retained body is delivered exactly as the article's lines are written, so no omission receipt of any kind accompanies this package. Citations: the retained excerpts carry no citation command at all, so no bibliography entry of the article is transcribed and no reference is redistributed. Reference adaptation: none was needed and none was made; every reference inside the delivered unit resolves inside it, so no unresolved reference or citation is produced. Printed slips kept literal: no printed word, symbol, hypothesis or number of the counted statement or of its proof was altered. Layout and class: the article's own class is \texttt{amsart} with packages including longtable, array, multicol, graphicx, hyperref, amsmath, amssymb, amsbsy, amsfonts, amsthm, latexsym, amsopn, amstext, amsxtra; the wrapper is the lane's pinned \texttt{amsart} editorial wrapper at 10pt on A4, which restates the theorem-like environments the retained text needs and adds the article's own macro definitions that the retained text needs (\texttt{\textbackslash F}, \texttt{\textbackslash Tr}), with extra wrapper packages (bm, bbm, dsfont, xspace, cancel, mathtools). The delivered numbering is the unit's own. Assets: no retained span contains a figure environment, a graphics-inclusion command, a PDF-page inclusion, a table or a TikZ picture; no image file is copied into this package, the package declares no asset dependency and no image file is redistributed with it. Licence: version arXiv:2608.10776v1 of this article is distributed under the Creative Commons Attribution 4.0 International licence, as recorded in the arXiv licence metadata for this exact version and shown on the abstract page https://arxiv.org/abs/2608.10776v1. A full rights notice is delivered at licenses/arxiv-2608.10776-CC-BY-4.0.txt. Characters: every retained excerpt is pure ASCII, so the delivered engine, XeLaTeX with its default Latin Modern text font, reports no missing character.
   \begin{thm}\label{T31}
    Let $q=2^m$, where $m$ is a positive integer. Let $f(X)=X+\gamma \Tr_{q}^{q^3}(X^2+X^{q+1}) \in \F_{q^3}[X].$ Then $f(X)$ is a polynomial over $\F_{q^3}$ if and only if $\gamma \in \F_q.$ 
    %Moreover, $f(X)$ is an involution when $\gamma \in \F_q$.
\end{thm}

\begin{proof}
    It is immediate that $f(X)$ is a permutation polynomial over $\F_{q^3}$ for $\gamma=0.$ We restrict to the case $\gamma \neq 0$ and aim to prove that $f(X)$ is a permutation polynomial if and only if $\gamma \in \F_q.$ For this, we will show that  for any $a \in \F_{q^3}$, the equation $$f(X)=a$$  has a unique solution in $\F_{q^3}$ if and only if $\gamma \in \F_q.$ Now, the equation $f(X)=a$ implies that
    \begin{equation}\label{e36}
    X+\gamma\Tr_{q}^{q^3}(X^2+X^{q+1})=a.    
    \end{equation}
    Let $\frac{X+a}{\gamma}=u:=\Tr_{q}^{q^3}(X^2+X^{q+1}) \in \F_q$. We substitute $X=u\gamma+a$ into Equation \eqref{e36}. This yields 
    \begin{equation*}
        u+\Tr_{q}^{q^3}(u^2\gamma^2+a^2+u^2\gamma^{q+1}+u\gamma^qa+u\gamma a^q+a^{q+1})=0,
    \end{equation*}
    which is equivalent to 
    \begin{equation}\label{e37}
        u^2\Tr_{q}^{q^3}(\gamma^2+\gamma^{q+1})+u(\Tr_{q}^{q^3}(\gamma^qa+\gamma a^q)+1)+\Tr_{q}^{q^3}(a^2+a^{q+1})=0.
    \end{equation}
   Consequently, $f(X)$ permutes $\F_{q^3}$ if and only if Equation \eqref{e37} has a unique solution in $\F_q$ for any $a \in \F_{q^3}$. 

If $\gamma \in \F_q$, then $u=\Tr_{q}^{q^3}(a^2+a^{q+1})$ is a unique solution of Equation \eqref{e37}. Hence, $f(X)$ is a permutation polynomial for $\gamma \in \F_q.$ We now assume that $\gamma \in \F_{q^3} \setminus \F_q.$

\textbf{Case 1.} If $\Tr_{q}^{q^3}(\gamma^2) \neq \Tr_{q}^{q^3}(\gamma^{q+1})$. Then for $a \in \F_q$, Equation \eqref{e37} has two distinct solutions, namely, $$u=0 \text{ and }  u=\frac{1}{\Tr_{q}^{q^3}(\gamma^2+\gamma^{q+1})}.$$ Therefore, $f(X)$ is not a permutation polynomial over $\F_{q^3}$.

\textbf{Case 2.} Suppose that $\Tr_{q}^{q^3}(\gamma^2)=\Tr_{q}^{q^3}(\gamma^{q+1})$.  The equality  $$\Tr_{q}^{q^3}(\gamma^2)=\Tr_{q}^{q^3}(\gamma^{q+1})$$ implies that 
\begin{equation*}
    \gamma^2(1+\gamma^{2q-2}+\gamma^{2q^2-2}+\gamma^{q-1}+\gamma^{q^2+q-2}+\gamma^{q^2-1})=0.
\end{equation*}
Since $\gamma \neq 0$, this reduces to 
$$1+\gamma^{2q-2}+\gamma^{2q^2-2}+\gamma^{q-1}+\gamma^{q^2+q-2}+\gamma^{q^2-1}=0.$$
 Let $y=\gamma^{q-1}$, then above equation becomes 
 \begin{equation*}
     1+y^2+y^{2(q+1)}+y+y^{q+2}+y^{q+1}=0.
 \end{equation*}
 We now use $z=y^{q+1}$ in the above equation to obtain
 \begin{equation}\label{e38}
     z^2+(y+1)z+y^2+y+1=0.
 \end{equation}
Equation \eqref{e38} has a solution in $\F_{q^3}$ if and only if $\Tr_{2}^{2^{3m}}\left(\frac{y^2+y+1}{y^2+1}\right)=0.$ Observe that 
$$\Tr_{2}^{2^{3m}}\left(\frac{y^2+y+1}{y^2+1}\right)=\Tr_{2}^{2^{3m}}\left(\frac{y}{y^2+1}+1\right).$$
When $m$ is odd, we have
$$\Tr_{2}^{2^{3m}}\left(\frac{y}{y^2+1}+1\right)=\Tr_{2}^{2^{3m}}\left(\frac{y}{y^2+1}\right)+1.$$
Moreover, $$\Tr_{2}^{2^{3m}}\left(\frac{y}{y^2+1}\right)+1=\Tr_{2}^{2^{3m}}\left(\frac{1}{y+1}+\frac{1}{y^2+1}\right)+1=1.$$
Therefore, if $m$ is odd, Equation \eqref{e38} has no solution in $\F_q$. Thus, there does not exist any $\gamma \in \F_{q^3} \setminus \F_q$ such that $\Tr_{q}^{q^3}(\gamma^2)=\Tr_{q}^{q^3}(\gamma^{q+1}).$ Consequently, we assume that $m$ is an even positive integer. Since $\Tr_{q}^{q^3}(\gamma^2)=\Tr_{q}^{q^3}(\gamma^{q+1})$, therefore, Equation \eqref{e37} reduces to
\begin{equation}\label{e39}
    u(\Tr_{q}^{q^3}(\gamma^qa+\gamma a^q)+1)+\Tr_{q}^{q^3}(a^2+a^{q+1})=0.
\end{equation}
Let $a=\frac{1}{\gamma^q+\gamma^{q^2}}.$
Substituting this value into $\Tr_{q}^{q^3}(a^2+a^{q+1})$, we obtain 
\begin{equation*}
\begin{split}
\Tr_{q}^{q^3}\left(a^2+a^{q+1}\right)=&\frac{1}{\gamma^{2q}+\gamma^{2q^2}}+\frac{1}{\gamma^{2q^2}+\gamma^{2}}+\frac{1}{\gamma^{2}+\gamma^{2q}}+\frac{1}{(\gamma^{q}+\gamma^{q^2})(\gamma+\gamma^{q^2})}+\frac{1}{(\gamma+\gamma^{q})(\gamma+\gamma^{q^2})}+\\&\frac{1}{(\gamma^{q}+\gamma^{q^2})(\gamma+\gamma^{q})},
\end{split}
\end{equation*}
which is equivalent to
\begin{equation}\label{e310}
\begin{split}
\Tr_{q}^{q^3}\left(a^2+a^{q+1}\right)=&
\frac{1}{\gamma^{2q}+\gamma^{2q^2}}+\frac{1}{\gamma^{2q^2}+\gamma^{2}}+\frac{1}{\gamma^{2}+\gamma^{2q}}+\frac{1}{\gamma^{q^2+q}+\gamma^{2q^2}+\gamma^{q+1}+\gamma^{q^2+1}}+\\&\frac{1}{\gamma^{q^2+1}+\gamma^{2}+\gamma^{q^2+q}+\gamma^{q+1}}+\frac{1}{\gamma^{q+1}+\gamma^{q^2+1}+\gamma^{2q}+\gamma^{q^2+q}}.
\end{split}
\end{equation}
As $\Tr_{q}^{q^3}(\gamma^2)=\Tr_{q}^{q^3}(\gamma^{q+1})$, it follows that 
\begin{equation}\label{e311}
    \gamma^2+\gamma^{2q}+\gamma^{2q^2}=\gamma^{q+1}+\gamma^{q^2+q}+\gamma^{q^2+1}.
\end{equation}
By using Equation \eqref{e311} into Equation \eqref{e310}, we conclude that
$$\Tr_{q}^{q^3}(a^2+a^{q+1})=0.$$
Moreover, for $a=\frac{1}{\gamma^q+\gamma^{q^2}}$, we have 
$$\Tr_{q}^{q^3}(\gamma^qa+\gamma a^q)+1=\Tr_{q}^{q^3}\left(a(\gamma^q+\gamma^{q^2})\right)+1=0.$$  Therefore, Equation \eqref{e39} admits $q$ distinct solutions in $\F_q$ for $a=\frac{1}{\gamma^q+\gamma^{q^2}}.$ Consequently, $$f(X)=\frac{1}{\gamma^q+\gamma^{q^2}}$$ has at least $q$ solutions in $\F_{q^3}$ and thus in this case also, $f(X)$ is not a permutation polynomial over $\F_{q^3}.$

%We now show that $f(X)$ is an involution when $\gamma \in \F_q.$ Let $u=\Tr_{q}^{q^3}(X^2+X^{q+1})$. Then 
%$$f(f(X))=X+\gamma u+\gamma \Tr_{q}^{q^3}(X^2+\gamma^2u^2+X^{q+1}+X^q\gamma u+X\gamma u+\gamma^2 u^2)$$
%$$=X+\gamma u+\gamma u+\gamma \Tr_{q}^{q^3}(X^q\gamma u+\gamma u X)$$
%$$=X$$
%as $\gamma, u \in \F_q.$ Hence, $f(X)$ is an involution over $\F_{q^3}.$
\end{proof}

\end{document}
